\begin{table}[ht]
\centering
\footnotesize
\caption{OMM Alignment with Becker et al. (2009) Design Requirements}
\label{tab:omm-becker-alignment}
\begin{tabularx}{\textwidth}{|l|X|}
\hline
\textbf{Design Requirement} & \textbf{OMM Instantiation in this Thesis} \\ \hline
1. Comparison with existing models & Our literature synthesis (Chapter~\ref{chap:related-work}) mapped existing siloed APM and monitoring maturity models, highlighting the lack of cloud-native and correlation-focused frameworks. \\ \hline
2. Iterative research strategy & The OMM was refined cyclically across the design and prototyping phases of the PANOPTES architecture. \\ \hline
3. Evaluation & The model's target state (Level~4) is evaluated through a controlled experiment using Chaos Engineering (Chapters~\ref{chap:experiment-implementation} and~\ref{chap:results-and-discussion}). \\ \hline
4. Multi-method presentation & Developed via a systematic literature synthesis (Rapid Review) and evaluated via quantitative systems engineering experiments. \\ \hline
5. Detailed documentation & All stages of the design, metric definitions, MADR records, and evaluation scripts are fully documented and public. \\ \hline
6. Target audience definition & The model explicitly targets Site Reliability Engineers (SREs), system architects, and operations managers. \\ \hline
7. Problem definition & Addresses the complexity and fragmentation of telemetry data ($O = \{M, L, T, E\}$) in ephemeral Kubernetes environments. \\ \hline
8. Design results & Delivers a prescriptive framework capable of guiding organizations from siloed monitoring to unified, correlated observability. \\ \hline
\end{tabularx}
\par\medskip
\ABNTEXfontereduzida\selectfont\textbf{Source:} Developed by author (2026)
\end{table}
